%==============================================================================
% Exemplo de Proposta de Trabalho de Conclusão de Curso com a classe texufpel
%
% A proposta tem o mesmo tema do exemplo de TCC (tcc/exemplo-tcc.tex): é o
% projeto que, depois de executado, resultaria naquele trabalho. Troque os
% dados do preâmbulo (os comentários "ALTERE" indicam o que mudar) e o texto
% dos capítulos.
%
% Compilação (a partir da pasta raiz do repositório):
%   pdflatex tcc/exemplo-tcc-proposta
%   bibtex   exemplo-tcc-proposta
%   pdflatex tcc/exemplo-tcc-proposta
%   pdflatex tcc/exemplo-tcc-proposta
%
% As páginas citadas nos comentários são as do Manual de normas UFPel para
% trabalhos acadêmicos (2023).
%==============================================================================

% Opção tcc-proposta: a proposta não é depositada na biblioteca, por isso a
% classe não gera capa (a opção capa é ignorada), nem referência do trabalho
% no resumo, e não há ficha catalográfica nem folha de aprovação.
\documentclass[tcc-proposta]{texufpel}

%------------------------------------------------------------------------------
% Dados institucionais. ALTERE para o seu curso.
% Unidade feminina (Faculdade, Escola): \unidade[a]{...}.
%------------------------------------------------------------------------------
\unidade{Centro de Desenvolvimento Tecnológico}
\curso{Ciência da Computação}
\nomecurso{Bacharelado em Ciência da Computação}
\titulocurso{Bacharel em Ciência da Computação}

%------------------------------------------------------------------------------
% Dados do trabalho. ALTERE.
% Os nomes abaixo são de grandes nomes da Computação e servem só de exemplo.
% Substitua pelos nomes da sua orientação.
%------------------------------------------------------------------------------
\title{Rotas de coleta seletiva com algoritmos de caminho mínimo: um estudo de caso em Pelotas}

\author{Lovelace}{Ada}                          % ALTERE: autor(a)
\advisor[Prof.~Dr.]{Dijkstra}{Edsger W.}        % ALTERE: orientador(a)
\coadvisor[Prof.ª~Dr.ª]{Hopper}{Grace}          % ALTERE: coorientador(a), opcional

\begin{document}

% Orientadora ou coorientadora: troque o rótulo da folha de rosto.
% \renewcommand{\advisorname}{Orientadora}
\renewcommand{\coadvisorname}{Coorientadora}

\maketitle   % folha de rosto (sem capa, por ser proposta)
\sloppy      % evita linhas que ultrapassam a margem direita

%==============================================================================
\chapter{Dados de Identificação}
%==============================================================================

\section{Nome do Projeto}
Rotas de coleta seletiva com algoritmos de caminho mínimo: um estudo de caso em
Pelotas.

\section{Local de Realização}
Centro de Desenvolvimento Tecnológico, Universidade Federal de Pelotas (UFPel).

\section{Responsável pelo Projeto}
Ada Lovelace

ada.lovelace@example.com % ALTERE: seu e-mail

\section{Professor Orientador}
Prof. Dr. Edsger W. Dijkstra

\section{Professora Coorientadora}
% Sem coorientação, apague esta seção.
Prof.ª Dr.ª Grace Hopper

%==============================================================================
\chapter{Sumário Executivo}
% No máximo uma página.
%==============================================================================

A coleta seletiva de Pelotas segue rotas definidas pela experiência das
equipes, sem apoio de ferramentas de otimização. Em bairros com muitas ruas de
mão única, a ordem em que o caminhão visita os pontos de coleta pode aumentar
bastante a distância rodada por dia. Este projeto propõe tratar o problema
como um problema de caminho mínimo em grafos: a malha viária de cada setor de
coleta será modelada como um grafo ponderado, e o algoritmo de
\citet{Dijkstra:1959:NTP} calculará a menor rota entre a garagem e a
cooperativa de reciclagem, passando pelos pontos do setor.

O objetivo principal é verificar se as rotas calculadas são mais curtas do que
as rotas praticadas hoje. A contribuição esperada é um protótipo em Python,
aplicado a quatro setores da cidade, e uma comparação, setor a setor, entre a
distância das rotas calculadas e a das rotas atuais. Se o ganho se confirmar,
o protótipo pode servir de ponto de partida para uma ferramenta de apoio ao
planejamento da coleta.

%==============================================================================
\chapter{Histórico e Justificativa}
% No mínimo uma página.
%==============================================================================

O problema do caminho mínimo é um dos mais antigos da teoria dos grafos
aplicada à computação. Em 1959, Dijkstra publicou, em uma nota de três
páginas, o algoritmo que leva seu nome, para encontrar o menor caminho entre
dois vértices de um grafo com arestas de comprimento conhecido
\citep{Dijkstra:1959:NTP}. No mesmo período surgiram outras soluções, como a
de Bellman, que admite pesos negativos \citep{Cormen:2022:IA}. Desde então, o
algoritmo de Dijkstra é tema obrigatório dos livros de algoritmos e base de
sistemas de navegação e de roteamento em redes de computadores
\citep{Knuth:1997:ACP,Cormen:2022:IA}.

A coleta seletiva é um caso em que esse problema aparece de forma direta. O
caminhão parte da garagem, passa por pontos de coleta e leva o material até a
cooperativa de reciclagem. Cada quilômetro a mais na rota é combustível,
tempo da equipe e desgaste do veículo. Em cidades de porte médio, como
Pelotas, é comum que as rotas sejam definidas pela experiência de quem
dirige, o que funciona bem em traçados simples, mas tende a gerar desvios em
áreas com muitas ruas de mão única, como o centro histórico.

A justificativa do projeto é, portanto, prática. Um algoritmo conhecido, de
implementação simples e de custo computacional baixo, pode indicar rotas
melhores sem exigir equipamentos nem programas comerciais. O trabalho também
tem valor acadêmico para o curso: aplica conteúdos das disciplinas de
estruturas de dados e de algoritmos a um problema real da cidade, na mesma
linha de outros trabalhos da UFPel que usam a Computação para apoiar decisões
ambientais, como o de \citet{Cruz:2015:FCS}.

O que distingue esta proposta de um uso direto de aplicativos de navegação é
o foco na coleta: os pontos de coleta, a garagem e a cooperativa são vértices
fixos, e a comparação será feita com as rotas que a equipe já usa. Assim, o
resultado mede o ganho real que uma ferramenta desse tipo traria para o
serviço.

%==============================================================================
\chapter{Objetivos e Metas}
%==============================================================================

O objetivo geral é verificar se o algoritmo de Dijkstra reduz a distância
percorrida nas rotas de coleta seletiva de Pelotas.

Os objetivos específicos são:

\begin{enumerate}
  \item modelar a malha viária de quatro setores de coleta como grafos
        ponderados;
  \item implementar o algoritmo de Dijkstra e validá-lo em um grafo pequeno,
        com resultado conhecido;
  \item calcular as rotas mínimas de cada setor;
  \item comparar as rotas calculadas com as rotas praticadas pela equipe de
        coleta.
\end{enumerate}

As metas, que permitem acompanhar a execução do projeto, são:

\begin{enumerate}
  \item ter os grafos dos quatro setores conferidos com a equipe de coleta até
        maio;
  \item ter o protótipo validado até julho, junto com a entrega da monografia
        intermediária;
  \item ter a comparação dos quatro setores concluída até setembro.
\end{enumerate}

Os objetivos são resultados do trabalho, e não atividades: as atividades que
levam a cada um deles estão no Capítulo~\ref{cap:atividades}. As duas listas
acima usam alíneas, como pede o Manual (p. 71): o texto que as antecede termina
em dois pontos, cada alínea termina em ponto e vírgula e a última, em ponto.

%==============================================================================
\chapter{Metodologia}
%==============================================================================

O projeto será desenvolvido em quatro etapas, uma para cada objetivo
específico.

Na primeira, os setores de coleta serão escolhidos com a equipe responsável,
procurando incluir traçados diferentes: o centro, com quarteirões regulares e
muitas ruas de mão única, e bairros com loteamentos mais recentes. Para cada
setor, a malha viária será modelada como um grafo ponderado, com a garagem, os
pontos de coleta e a cooperativa como vértices e os trechos de rua como
arestas, com a distância em quilômetros como peso. O sentido das vias será
representado por arestas dirigidas.

Na segunda etapa, o algoritmo de Dijkstra será implementado em Python, com a
fila de prioridade do módulo \texttt{heapq}, seguindo a descrição de
\citet{Cormen:2022:IA}. A implementação será validada em um grafo pequeno,
cujas distâncias mínimas podem ser calculadas à mão.

Na terceira, o protótipo calculará a rota mínima de cada setor. Na quarta, as
rotas calculadas serão comparadas com as rotas praticadas, informadas pela
equipe de coleta, pela distância total percorrida em cada setor.

%==============================================================================
\chapter{Plano de Atividades e Cronograma}
\label{cap:atividades}
% Para um ano.
%==============================================================================

As atividades previstas são:

\begin{enumerate}
  \item revisão da literatura sobre caminho mínimo e roteamento;
  \item levantamento dos setores e das rotas praticadas, com a equipe de
        coleta;
  \item modelagem dos grafos dos setores;
  \item implementação e validação do algoritmo de Dijkstra;
  \item cálculo das rotas e comparação com as rotas praticadas;
  \item escrita da monografia;
  \item entrega da monografia intermediária;
  \item entrega da monografia final;
  \item apresentação final (banca);
  \item submissão de artigo a evento de iniciação científica (opcional).
\end{enumerate}

A Tabela~\ref{tab:cronograma} distribui as atividades ao longo de 2026,
identificadas pela letra da alínea. As
atividades de escrita, entrega e apresentação da monografia são obrigatórias,
e os meses de entrega e da banca seguem o calendário definido pelo professor
responsável pelo Trabalho de Conclusão de Curso.

% Cronograma em tabela: uma coluna para a atividade e uma para cada mês.
% A coluna da atividade usa só a letra da alínea da lista acima, para caber
% na página.
% Os meses vão abreviados conforme o Anexo A do Manual (p. 79).
\begin{table}[htbp]
  \caption{Cronograma de atividades em 2026}
  \label{tab:cronograma}
  \centering
  \setlength{\tabcolsep}{4pt}
  \begin{tabular}{l*{12}{c}}
    \hline
    Atividade & jan. & fev. & mar. & abr. & maio & jun. & jul. & ago. & set. & out. & nov. & dez. \\
    \hline
    a) & X & X & X &   &   &   &   &   &   &   &   &   \\
    b) &   & X & X & X &   &   &   &   &   &   &   &   \\
    c) &   &   &   & X & X &   &   &   &   &   &   &   \\
    d) &   &   &   &   & X & X & X &   &   &   &   &   \\
    e) &   &   &   &   &   &   &   & X & X &   &   &   \\
    f) &   &   & X & X & X & X & X & X & X & X & X &   \\
    g) &   &   &   &   &   &   & X &   &   &   &   &   \\
    h) &   &   &   &   &   &   &   &   &   &   & X &   \\
    i) &   &   &   &   &   &   &   &   &   &   &   & X \\
    j) &  &   &   &   &   &   &   &   &   & X &   &   \\
    \hline
  \end{tabular}
  \caption*{Fonte: Elaborada pela autora (2026).}
\end{table}

A tabela segue o padrão do Instituto Brasileiro de Geografia e Estatística
indicado pelo Manual (p. 17): título em cima, fonte embaixo, traços
horizontais no cabeçalho e no fim, sem traços verticais. A especificação
\verb|l*{12}{c}| cria uma coluna à esquerda e doze colunas centralizadas, uma
por mês.

%------------------------------------------------------------------------------
% Bibliografia: estilo abnt.bst e base bibliografia.bib, na raiz. Só as obras
% citadas no texto entram na lista. O título da seção segue o Manual, que usa
% a expressão "Referências" (p. 21).
%------------------------------------------------------------------------------
\bibliographystyle{abnt}
\bibliography{bibliografia}

%------------------------------------------------------------------------------
% Assinaturas da proponente e do orientador.
%------------------------------------------------------------------------------
\chapter*{Assinaturas}

\vspace{2cm}

\begin{center}
\rule{8cm}{.3mm}

Ada Lovelace\\          % ALTERE
Proponente
\end{center}

\vspace{3cm}

\begin{center}
\rule{8cm}{.3mm}

Prof. Dr. Edsger W. Dijkstra\\   % ALTERE
Orientador
\end{center}

\end{document}
